\documentclass[11pt,a4paper]{article}
\usepackage[T1]{fontenc}
\usepackage[english]{babel}
\usepackage{csquotes}
\usepackage{amsmath}
\usepackage{isabelle,isabellesym}

\usepackage[top=3cm,bottom=4.5cm]{geometry}

% this should be the last package used
\usepackage{pdfsetup}

\urlstyle{rm}
\isabellestyle{tt}

\title{Byzantine-tolerant detection of causality:\\
       there is no holy grail}
\author{Mechanisation in Isabelle/HOL of the impossibility results of\\
        Misra and Kshemkalyani (Parallel Computing 124, 2025)}
\date{\today}

\begin{document}

\maketitle

\begin{abstract}
We mechanise in Isabelle/HOL the impossibility results of
Misra and Kshemkalyani, \emph{Byzantine-tolerant detection of
causality: there is no holy grail} (Parallel Computing 124 (2025),
article 103136).  The causality determination problem
$\mathrm{CD}(E,F,e_i^*)$ asks an asynchronous distributed algorithm
to reconstruct the global execution history $E$ at a correct
process $p_i$, with respect to a target event $e_i^*$, such that
the happened-before relation evaluated against the algorithm's
collected $F$ agrees with the truth in $E$.  In the presence of
even a single Byzantine process this is impossible.

We formalise:

\begin{itemize}
  \item Theorems 1 and 2 (paper Section 4.1): no algorithm can
        prevent false negatives in the Byzantine model, and for
        internal events no algorithm can prevent both false
        negatives and false positives.  Both proofs are
        constructive: given any candidate algorithm we exhibit a
        Byzantine adversary using a fresh natural number as event
        identifier.
  \item Theorems 3, 4, 5 (paper Section 4.2): the headline
        impossibility results for unicast, broadcast, and
        multicast communication.  The proofs follow the paper's
        chain $\mathrm{Consensus} \preceq \mathrm{Black\_Box} \preceq
        \mathrm{CD}$, with the FLP impossibility imported from the
        AFP entry \emph{FLP} (Bisping et al., 2025) as a proven
        theorem.
\end{itemize}

The development sits on top of the AFP entry \emph{FLP}.  Two
named meta-level hypotheses are exposed on the headline theorems:
\verb|cd_can_identify_correct| capturing the paper's argument that
solving CD requires identifying the correct set, and
\verb|bb_realizes_flp_consensus| capturing the standard reduction
that an abstract Black\_Box solver can be realised by an
asynchronous distributed protocol.  Neither is an internal HOL
axiom; both are faithful to the paper's prose and discharged at
interpretation time by the user.

\medskip

\noindent\emph{License:} BSD-3-Clause (matching the AFP \emph{FLP}
entry).  The source paper is open access under CC-BY 4.0 and
redistributed in this repository under that license.
\end{abstract}

\tableofcontents

\newpage

\section*{Reading order}

The theories are intended to be read in the following order:

\begin{enumerate}
  \item \verb|ByzantineSystem|: the process partition (correct
        $\uplus$ Byzantine) and the abstract Consensus signature.
  \item \verb|Events|: events, per-process and global histories,
        program order, message order, happened-before relation
        (paper Definitions 1, 2).
  \item \verb|CD|: the Causality Determination problem,
        $\mathit{valid}(F)$, false positives / negatives, the
        adversary model, the CD-solver signature (paper
        Definition 5).
  \item \verb|BlackBox|: the Black\_Box problem of paper Section
        4.2 with its piecewise broadcast value $w$.
  \item \verb|Reductions|: the two reductions of paper Section
        4.2 -- R1 (Consensus $\preceq$ BB) constructive in
        declarative Isar, R2 (BB $\preceq$ CD) constructive modulo
        the paper's named meta-level step
        \verb|cd_can_identify_correct|.
  \item \verb|Theorems_1_2|: paper Section 4.1, Theorems 1 and 2,
        with explicit Byzantine adversary constructions using
        fresh natural numbers.
  \item \verb|FLP_Consensus|: the FLP-style consensus predicate
        with the impossibility proven against the AFP entry's
        \verb|ConsensusFails|, and the bridge predicate
        \verb|bb_realizes_flp_consensus|.
  \item \verb|Impossibility|: paper Section 4.2, Theorems 3, 4,
        5 -- the headline results.
  \item \verb|Foundation_Vacuity|: a regression diagnostic
        documenting why the abstract \verb|solves_Consensus|
        predicate is too weak to express FLP, and motivating the
        FLP-style chain used by \verb|Impossibility|.
\end{enumerate}

\newpage

% include generated text of all theories
\input{session}

\nocite{*}
\bibliographystyle{abbrv}
\bibliography{root}

\end{document}
